\section*{Abstract}

Remote train operation is becoming a critical capability as railways move toward higher grades of automation, and at those levels a reliable low latency video feed is the operator's primary source of situational awareness. The remote operation project at NTNU previously depended on a closed commercial streaming product that prevented configuration control and internal latency measurement, while existing research offered little guidance on which streaming configurations suit railway use. This thesis presents the design, implementation and experimental evaluation of a custom GStreamer based video streaming pipeline built to replace that product. The pipeline streams three cameras from a sender computer over a 5G connection to a remote receiving computer, and a measurement framework based on GStreamer probes and RTP timestamps captures sender, transit and receiver latency separately, bringing the measurement uncertainty down.

The system was evaluated across a wide range of configurations, covering the H.264, H.265 and MJPEG codecs, CPU and GPU encoding, 1080p and 720p resolution, transcoded and passthrough operation, the UDP and TCP transport protocols, SRTP encryption and constrained bandwidth. Each configuration was assessed on end-to-end latency, packet loss, throughput, sender CPU usage and visual quality. Low latency streaming was proved achievable across nearly every configuration. Among the transcoded configurations H.265 with GPU acceleration stood out with the lowest end-to-end latency. Under normal conditions UDP delivered lower latency than TCP while TCP produced cleaner images, but under constrained bandwidth UDP degraded quality whereas TCP fell behind. Adding SRTP encryption carried almost no cost in latency or CPU usage.

The recommended configuration for remote train operation is to use internal camera encoding with H.265 codec over UDP transport layer. If it is required for the computer to encode or re-encode, H.265 encoding with GPU acceleration over UDP transport layer is recommended TCP should be reserved for stable high bandwidth conditions only. With the constant camera capture and display delays included, the best configurations reach a glass-to-glass latency of around 80 to 100~ms, below the 120 to 240~ms reported by earlier work in the project. The resulting pipeline gives the project full control over every stage of the stream and a stable foundation for future work, though confirming whether the achieved latency and image quality are sufficient for safe driving still requires a human factors evaluation with real operators.

%cspell:disable

----------------------- Norwegian~-----------------------

Når vi nærmer oss høyere grad av automatisering i jernbaneindustrien blir fjernstyring av tog viktigere. A ha en pålitelig og rask videostrømning som operatørene kan stole på og bruke er helt nødvendig. Fjernstyringsprosjektet på NTNU har tidligere vært avhengig av kommersielle operatører, noe som har bydd på utfordringer rundt målinger og resultater. Ikke minst for å justere operasjonen til å tilpasse seg fjernstyring av spesifikt tog. Denne masteroppgaven tar for seg designet, implementasjonen og vurderingen av en GStreamer basert videostrømningsystem lagd for å erstatte tidligere brukte produkter. Denne løsningen strømmer tre videoer fra forskjellige kameraer fra en PC over 5G til en annen PC og har innebygde målingsinstrumenter for å finne en nøyaktig og oppstykket forsinkelse av videoen.

Systemet var evaluert på tvers av mange forskjellige konfigurasjoner og dekker enkodere som H.264, H.265 og MJPEG, CPU og GPU enkoding, 1080p og 720p oppløsning, omkoding og uten ekstra omkoding, UDP og TCP nettverk transport protokoller, SRTP kryptering og begrenset båndbredde. Hver konfigurasjon var vurdert basert på ende til ende forsinkelse, nettverks pakketap, datagjennomstrømning, CPU bruk på avsender og visuell kvalitet. Det var mulig å oppnå relativ lav forsinkelse av videostrømningen på nesten alle konfigurasjonene. Av omkoding konfigurasjonene så var H.265 med hjelp av GPU veldig vellykket, og UDP resulterte i lavere forsinkelse enn TCP, mens TCP produserte bedre bildekvalitet. Under nettverksbegrensninger så falt kvaliteten med UDP mens TCP ble forsinket. Å legge til kryptering påvirket veldig lite. 

Den anbefalte løsningen for fjernstyring av tog er sending av kamera kodet H.265 videostrøm over UDP, eller H.265 med GPU hvis det er nødvendig med omkoding. Hvis enormt stabilt bredbånd er tilgjengelig kan bruk av TCP også utforskes. Med kamera bildeopptak og skjerm visning inkludert i treghetsmålinger av videostrømmen er det mulig å få til en gjennomsnittelig 80~ms forsinkelse, noe som er under kravet 100~ms fra 3GPP TS 22.289. Det er også lavere forsinkelse enn tidligere målt i prosjektet utført av NTNU. Den foreslåtte konfigurasjonen av systemet gir prosjektet full kontroll over hvert stadie av videostrømmingen og danner grunnlag for framtidig arbeid. 